%! AP Statistics — Unit 9, Lesson 9.5 — Exit Ticket (Answer Key)
\documentclass[10pt]{article}
\usepackage{apstats-article}
\usepackage{apstats-key}

\begin{document}

\pageheader{Unit 9, Lesson 9.5}{Exit Ticket --- Answer Key}
\namedateperiod

\vspace{0.1in}

A meteorologist fits a regression line relating elevation ($x$, in hundreds of feet)
to average annual temperature ($y$, ${}^\circ$F) for a random sample of 18 weather
stations. Assume LINER conditions are met. Computer output:

\begin{center}
\renewcommand{\arraystretch}{1.4}
\begin{tabular}{lrrrr}
    \textbf{Predictor} & \textbf{Coef} & \textbf{SE Coef} & \textbf{T} & \textbf{P} \\
    \hline
    Constant   & 65.4 & 3.2  & 20.44 & 0.000 \\
    Elevation  & $-2.1$ & 0.45 & $-4.67$ & 0.000 \\
\end{tabular}
\end{center}

\begin{enumerate}

  \item State H$_0$ and H$_a$ for a test of the true slope of the regression line
        relating temperature to elevation.

        H$_0$: \ans{$\beta = 0$; there is no linear relationship between elevation and
        average annual temperature for weather stations.}

        \vspace{0.05in}

        H$_a$: \ans{$\beta \ne 0$; there is a linear relationship between elevation and
        average annual temperature for weather stations.}

        \vspace{0.1in}

  \item Interpret the $p$-value ($= 0.000$) in context, assuming H$_0$ is true.

        \ans{Assuming the true slope is 0 (no linear relationship between elevation and
        average annual temperature), there is approximately a 0.000 probability of
        obtaining a sample slope as extreme as $b = -2.1\,{}^\circ$F per hundred feet
        just by chance.}

  \item State a conclusion at $\alpha = 0.05$.

        \ans{Since $p \approx 0 < 0.05 = \alpha$, we reject H$_0$. There is convincing
        evidence that the true slope of the population regression line relating average
        annual temperature to elevation for weather stations is not zero.}

\end{enumerate}

\begin{teachernote}
$p \approx 0.000$ means the $p$-value rounds to 0.000 at three decimal places, not that
it is exactly zero. Students may write $p < 0.001$ or $p \approx 0.000$ --- either phrasing
earns full credit. Remind students that the $t$-statistic here is $-4.67$ with $df = 16$.
\end{teachernote}

\end{document}
